% Selected original problems and all proof adaptations:
% authoring/notes/unit04-exercise-sources.md.
All vector spaces below are finite-dimensional real spaces.

\subsection*{Linear maps and their kernels}

% Lang LA III.2.1(a-g); calculus part (h) not selected.
\begin{problem}[Testing linearity]\label{ex:04:linearity}
Determine which maps are linear. Give a proof or a counterexample in each case.
\begin{enumerate}
\item \(F:\R^3\to\R^2,\quad F(x,y,z)=(x,z)^\top\).
\item \(F:\R^4\to\R^4,\quad F(x)=-x\).
\item \(F:\R^3\to\R^3,\quad F(x)=x+(0,-1,0)^\top\).
\item \(F:\R^2\to\R^2,\quad F(x,y)=(2x+y,y)^\top\).
\item \(F:\R^2\to\R^2,\quad F(x,y)=(2x,y-x)^\top\).
\item \(F:\R^2\to\R^2,\quad F(x,y)=(y,x)^\top\).
\item \(F:\R^2\to\R,\quad F(x,y)=xy\).
\end{enumerate}
\end{problem}
\begin{solution}
The maps in (a), (b), (d), (e), and (f) are respectively multiplication by
\[
\begin{pmatrix}1&0&0\\0&0&1\end{pmatrix},\quad
-I_4,\quad
\begin{pmatrix}2&1\\0&1\end{pmatrix},\quad
\begin{pmatrix}2&0\\-1&1\end{pmatrix},\quad
\begin{pmatrix}0&1\\1&0\end{pmatrix}.
\]
Each is linear because \(A(su+tv)=sAu+tAv\).
In (c), \(F(0)=(0,-1,0)^\top\ne0\), which rules out linearity.
In (g), \(F(1,0)=F(0,1)=0\) but \(F(1,1)=1\), so additivity fails.
\end{solution}

% Lang LA III.2.18(a-c) and III.2.19(a-c), original data.
\begin{problem}[Recovering a map from two values]\label{ex:04:two-values}
In each case \(L:\R^2\to\R^2\) is linear. Find both \(L(1,0)\) and \(L(0,1)\).
\begin{enumerate}
\item \(L(3,1)=(1,2)^\top,\quad L(-1,0)=(1,1)^\top\).
\item \(L(4,1)=(1,1)^\top,\quad L(1,1)=(3,-2)^\top\).
\item \(L(1,1)=(2,1)^\top,\quad L(-1,1)=(6,3)^\top\).
\end{enumerate}
\end{problem}
\begin{solution}
Express the standard columns using the given input columns, then apply \(L\).
\begin{enumerate}
\item The identities
\[
e_1=-(-1,0)^\top,\qquad e_2=(3,1)^\top+3(-1,0)^\top
\]
give \(L(e_1)=(-1,-1)^\top\) and \(L(e_2)=(4,5)^\top\).
They check the data: \(3L(e_1)+L(e_2)=(1,2)^\top\) and
\(-L(e_1)=(1,1)^\top\).
\item Here
\[
e_1=\tfrac13(4,1)^\top-\tfrac13(1,1)^\top,\qquad
e_2=-\tfrac13(4,1)^\top+\tfrac43(1,1)^\top.
\]
Thus \(L(e_1)=(-2/3,1)^\top\), \(L(e_2)=(11/3,-3)^\top\).
The combinations \(4L(e_1)+L(e_2)\) and \(L(e_1)+L(e_2)\)
are the two prescribed outputs.
\item Finally,
\[
e_1=\tfrac12(1,1)^\top-\tfrac12(-1,1)^\top,\qquad
e_2=\tfrac12(1,1)^\top+\tfrac12(-1,1)^\top.
\]
Hence \(L(e_1)=(-2,-1)^\top\), \(L(e_2)=(4,2)^\top\).
Their sum and their difference \(L(e_2)-L(e_1)\) recover the given outputs.
\end{enumerate}
\end{solution}

% Lang LA III.2.6, including the requested generalization to finitely many coordinates.
\begin{problem}[Combining scalar-valued maps]\label{ex:04:coordinate-maps}
Let \(f,g:V\to\R\) be linear. Prove that
\(F:V\to\R^2\), \(F(v)=(f(v),g(v))^\top\), is linear.
Generalize to any finite number of scalar-valued linear maps.
\end{problem}
\begin{solution}
For linear \(f_1,\ldots,f_p:V\to\R\), \(p\geqslant1\), define
\(F(v)=(f_1(v),\ldots,f_p(v))^\top\).
For \(u,v\in V\) and \(a,b\in\R\), coordinate \(j\) gives
\[
f_j(au+bv)=af_j(u)+bf_j(v)\qquad(1\leqslant j\leqslant p).
\]
Equality in every coordinate proves \(F(au+bv)=aF(u)+bF(v)\).
The requested pair is the case \(p=2\).
\end{solution}

% Lang LA III.2.15.
\begin{problem}[Independent images]\label{ex:04:independent-preimages}
Let \(F:V\to W\) be linear. Suppose \(w_1,\ldots,w_k\) are independent
and \(F(v_i)=w_i\) for every \(i\). Prove that \(v_1,\ldots,v_k\) are independent.
\end{problem}
\begin{solution}
If \(\sum_i a_iv_i=0\), applying \(F\) gives
\(\sum_i a_iw_i=0\). Independence of the \(w_i\) gives \(a_i=0\)
for every \(i\). This proves the required independence without assuming
that \(F\) is injective on all of \(V\).
\end{solution}

% Lang LA III.2.16 and III.2.17, both complete.
\begin{problem}[A scalar-valued map]\label{ex:04:scalar-kernel}
Let \(F:V\to\R\) be linear, let \(W=\ker F\ne V\), and fix \(v_0\notin W\).
\begin{enumerate}
\item Show that every \(v\in V\) can be written \(v=w+cv_0\), where
\(w\in W\) and \(c\in\R\).
\item Prove that \(W\) is a subspace. If \((v_1,\ldots,v_r)\) is a
basis of \(W\), prove that \((v_0,v_1,\ldots,v_r)\) is a basis of \(V\).
\end{enumerate}
\end{problem}
\begin{solution}
\begin{enumerate}
\item Since \(F(v_0)\ne0\), set
\[
c=\frac{F(v)}{F(v_0)},\qquad w=v-cv_0.
\]
Then \(F(w)=F(v)-cF(v_0)=0\), so \(w\in W\) and the formula holds.
\item We have \(F(0)=0\), and for \(u,w\in W\),
\(F(au+bw)=aF(u)+bF(w)=0\); hence \(W\) is a subspace.
Part (a) and the basis of \(W\) prove spanning.
In a relation \(a_0v_0+\sum_{i=1}^r a_iv_i=0\), applying \(F\) gives
\(a_0F(v_0)=0\), hence \(a_0=0\).
Independence in \(W\) gives all remaining coefficients zero.
The argument also covers \(W=\{0\}\), when its basis is empty.
\end{enumerate}
\end{solution}

% Harville 22.1.
\begin{problem}[Extending a linear map]\label{ex:04:extension}
Let \(U\subseteq V\) be a subspace and let \(S:U\to W\) be linear.
Show that there is a linear \(T:V\to W\) whose restriction to \(U\) is \(S\).
\end{problem}
\begin{solution}
\begin{itemize}
\item Choose a basis \((u_1,\ldots,u_r)\) of \(U\) and extend it to
\((u_1,\ldots,u_r,v_1,\ldots,v_s)\) of \(V\).
\item Choose any \(w_1,\ldots,w_s\in W\), and define
\[
T\left(\sum_{i=1}^r a_iu_i+\sum_{j=1}^s b_jv_j\right)
=\sum_{i=1}^r a_iS(u_i)+\sum_{j=1}^s b_jw_j.
\]
Unique basis coefficients make this well-defined; addition and scalar
multiplication of the coefficients prove linearity.
\item For \(u\in U\), all \(b_j=0\), so \(T(u)=S(u)\).
If \(U=V\) there are no new images to choose; if \(U=\{0\}\), the
first sum is empty. The new images \(w_j\) can be chosen freely.
\end{itemize}
\end{solution}

% Harville 22.3(a,b); retain the original direct spanning proof.
\begin{problem}[A subspace avoiding the kernel]\label{ex:04:kernel-complement}
Let \(T:V\to W\) be linear and \(U\subseteq V\) a subspace satisfying
\(U\cap\ker T=\{0\}\). Let \(u_1,\ldots,u_r\in U\) be independent.
\begin{enumerate}
\item Prove that \(T(u_1),\ldots,T(u_r)\) are independent.
\item If also \(r=\dim U\) and \(U\oplus\ker T=V\), prove that these
images form a basis of \(\operatorname{im}T\).
\end{enumerate}
\end{problem}
\begin{solution}
\begin{enumerate}
\item A relation \(\sum_i a_iT(u_i)=0\) gives
\[
\sum_i a_iu_i\in U\cap\ker T=\{0\}.
\]
Independence of the \(u_i\) forces every \(a_i=0\).
\item The \(u_i\) form a basis of \(U\).
For \(v\in V\), write \(v=u+z\) with \(u\in U\), \(z\in\ker T\),
and then \(u=\sum_i c_iu_i\).
Thus \(T(v)=T(u)=\sum_i c_iT(u_i)\), proving spanning.
Part (a) proves independence, including the empty list if \(T=0\).
\end{enumerate}
\end{solution}

% Harville 22.6.
\begin{problem}[Factoring through another map]\label{ex:04:factor-map}
Let \(T:V\to W\) and \(R:U\to W\) be linear.
If \(\operatorname{im}T\subseteq\operatorname{im}R\), prove that
there is a linear \(S:V\to U\) such that \(T=RS\).
\end{problem}
\begin{solution}
\begin{itemize}
\item Choose a basis \((v_1,\ldots,v_n)\) of \(V\).
For each \(i\), choose \(u_i\in U\) with \(R(u_i)=T(v_i)\),
which is possible by the inclusion of images.
\item Define \(S(\sum_i a_iv_i)=\sum_i a_iu_i\).
Unique coefficients make \(S\) well-defined; their addition and scaling
give linearity. For every \(v=\sum_i a_iv_i\),
\[
RS(v)=\sum_i a_iR(u_i)=\sum_i a_iT(v_i)=T(v).
\]
Thus \(T=RS\). If \(V=\{0\}\), the zero map \(S\) gives the same conclusion.
\end{itemize}
\end{solution}

% Lang LA III.4.8(a,b); original answer is only a kernel assertion.
\begin{problem}[Two invertible maps]\label{ex:04:invertible-maps}
Show that each of the following maps \(L:\R^3\to\R^3\) is invertible.
\begin{enumerate}
\item \(L(x,y,z)=(x-2y,x+z,x+y+2z)^\top\).
\item \(L(x,y,z)=(2x-y+z,x+y,3x+y+z)^\top\).
\end{enumerate}
\end{problem}
\begin{solution}
Both maps are linear by their coordinate formulas.
Solve \(L(x,y,z)=(u,v,w)^\top\) for arbitrary \(u,v,w\).
\begin{enumerate}
\item From \(x=u+2y\) and \(z=v-x\), the last equation becomes
\[
w=x+y+2(v-x)=2v-u-y.
\]
Hence
\[
y=2v-u-w,\qquad x=4v-u-2w,\qquad z=u-3v+2w.
\]
These equations determine the only possible preimage. Substitution into
the three coordinates gives \(u,v,w\), so it exists for every output.
\item The second equation gives \(y=v-x\), and the first gives
\(z=u+v-3x\). The last equation is then \(w=u+2v-x\).
Thus the unique candidate is
\[
x=u+2v-w,\qquad y=-u-v+w,\qquad z=-2u-5v+3w.
\]
Substitution gives \(2x-y+z=u\), \(x+y=v\), and \(3x+y+z=w\).
Again every output has exactly one preimage.
\end{enumerate}
\end{solution}

% Lang LA III.4.1; ordered basis retained.
\begin{problem}[A nonzero map whose square is zero]\label{ex:04:square-zero-map}
Let \(L:\R^2\to\R^2\) be linear, \(L\ne0\), and \(L^2=L\circ L=0\).
Show that there is a basis \((a,b)\) of \(\R^2\) such that
\(L(a)=b\) and \(L(b)=0\).
\end{problem}
\begin{solution}
Since \(L\ne0\), choose \(a\) with \(b=L(a)\ne0\).
Then \(L(b)=L^2(a)=0\).
If \(sa+tb=0\), applying \(L\) gives \(sb=0\), hence \(s=0\).
The original relation becomes \(tb=0\), so \(t=0\).
These two independent columns form the required basis of \(\R^2\).
\end{solution}

% Lang LA III.4.10; direct decomposition, without orthogonality.
\begin{problem}[An idempotent map]\label{ex:04:idempotent}
Let \(P:V\to V\) be linear and satisfy \(P^2=P\).
Prove
\[
V=\ker P+\operatorname{im}P,\qquad
\ker P\cap\operatorname{im}P=\{0\}.
\]
\end{problem}
\begin{solution}
For \(v\in V\),
\[
v=(v-Pv)+Pv,\qquad P(v-Pv)=Pv-P^2v=0.
\]
The first term lies in \(\ker P\), the second in \(\operatorname{im}P\);
therefore their sum is all of \(V\).
If \(w\) belongs to the intersection, write \(w=Pv\). Then
\[
w=Pv=P^2v=P(w)=0.
\]
This proves that the sum is direct.
\end{solution}

% Harville 22.7(a,b), including its general (not necessarily linear) maps.
\begin{problem}[A left inverse and a right inverse]\label{ex:04:two-inverses}
Let \(T:V\to W\), \(R:W\to V\), and \(S:W\to V\) be maps satisfying
\(RT=I_V\) and \(TS=I_W\).
\begin{enumerate}
\item Show that \(T\) is invertible.
\item Show that \(R=S=T^{-1}\).
\end{enumerate}
\end{problem}
\begin{solution}
\begin{enumerate}
\item If \(T(v)=T(v')\), applying \(R\) gives \(v=v'\).
For \(w\in W\), the element \(S(w)\) satisfies \(T(S(w))=w\).
Thus \(T\) is both injective and surjective.
\item For every \(w\in W\),
\[
R(w)=R(T(S(w)))=S(w).
\]
The identities \(RT=I_V\), \(TS=I_W\) therefore make this common map
the inverse of \(T\).
\end{enumerate}
The proof requires no linearity assumption on these maps.
\end{solution}

\subsection*{Coordinates and matrix representations}

% Lang LA IV.2.1(a,f), and IV.3.2; geometric terminology omitted.
\begin{problem}[Reading a representation from basis images]\label{ex:04:representation-columns}
Find the indicated representation matrices.
\begin{enumerate}
\item \(F:\R^4\to\R^2\), \(F(x_1,x_2,x_3,x_4)=(x_1,x_2)^\top\),
using standard bases.
\item \(F:\R^4\to\R^4\), \(F(x_1,x_2,x_3,x_4)=(x_1,x_2,0,0)^\top\),
using standard bases.
\item A linear map \(L:V\to V\), with a basis \(B=(v_1,\ldots,v_n)\) and
\(L(v_i)=c_iv_i\), using \(B\) for both domain and codomain.
\end{enumerate}
\end{problem}
\begin{solution}
Column \(j\) is the coordinate column of the image of input basis vector \(j\).
Thus the matrices are respectively
\[
\begin{pmatrix}1&0&0&0\\0&1&0&0\end{pmatrix},\qquad
\begin{pmatrix}1&0&0&0\\0&1&0&0\\0&0&0&0\\0&0&0&0\end{pmatrix},\qquad
\operatorname{diag}(c_1,\ldots,c_n).
\]
For (c), \([L(v_j)]_B=c_je_j\), which checks every diagonal and off-diagonal entry.
The first two maps have different codomains and hence different numbers of matrix rows.
\end{solution}

% Lang LA IV.3.1(b), full original data; (a) is a classroom example.
\begin{problem}[Changing coordinates]\label{ex:04:change-coordinates}
In \(\R^3\), let
\[
B=\bigl((3,2,1)^\top,(0,-2,5)^\top,(1,1,2)^\top\bigr),
\quad
C=\bigl((1,1,0)^\top,(-1,2,4)^\top,(2,-1,1)^\top\bigr).
\]
Find \([I]_{C\leftarrow B}\), the matrix converting \(B\)-coordinates
to \(C\)-coordinates.
\end{problem}
\begin{solution}
Let \(U,V\) have the columns of \(B,C\), respectively:
\[
U=\begin{pmatrix}3&0&1\\2&-2&1\\1&5&2\end{pmatrix},\qquad
V=\begin{pmatrix}1&-1&2\\1&2&-1\\0&4&1\end{pmatrix}.
\]
If \(U(a,b,c)^\top=0\), the first two rows give \(c=-3a\) and \(b=-a/2\);
the third then gives \(-15a/2=0\). Thus all three coefficients vanish,
so \(B\) is a basis.
To express \((r,s,t)^\top\) in \(C\), solve
\[
a-b+2c=r,\qquad a+2b-c=s,\qquad4b+c=t.
\]
Subtracting the first equation from the second and using the third gives
\[
b=\frac{s-r+3t}{15},\qquad
c=\frac{4r-4s+3t}{15},\qquad
a=\frac{2r+3s-t}{5}.
\]
The unique solution for every \((r,s,t)\) also proves that \(C\) is a basis.
Apply this formula to the three columns of \(U\):
\[
[I]_{C\leftarrow B}
=P=\begin{pmatrix}
11/5&-11/5&3/5\\
2/15&13/15&2/5\\
7/15&23/15&2/5
\end{pmatrix}.
\]
The product \(VP=U\) verifies all three coordinate expansions.
For example its first column is
\((11/5)(1,1,0)^\top+(2/15)(-1,2,4)^\top
+(7/15)(2,-1,1)^\top=(3,2,1)^\top\).
\end{solution}

% Harville 22.12.
\begin{problem}[Constructing a desired coordinate change]\label{ex:04:prescribed-change}
Let \(B=(v_1,\ldots,v_n)\) be a basis of \(V\), and let
\(A\in\R^{n\times n}\) be invertible. Write \(A^{-1}=(f_{ij})\) and set
\[
w_j=\sum_{i=1}^n f_{ij}v_i\qquad(1\leqslant j\leqslant n).
\]
Prove that \(C=(w_1,\ldots,w_n)\) is a basis of \(V\) and
\([I]_{C\leftarrow B}=A\).
\end{problem}
\begin{solution}
A relation \(\sum_j c_jw_j=0\) has \(B\)-coordinate column \(A^{-1}c\).
Thus \(A^{-1}c=0\), so \(c=0\).
The \(n\) independent vectors form a basis of the \(n\)-dimensional space.
For any \(v\in V\), expanding its \(C\)-coordinates in the \(v_i\) gives
\[
[v]_B=A^{-1}[v]_C,\qquad [v]_C=A[v]_B.
\]
This is exactly the claimed representation of the identity.
\end{solution}

% Harville 22.13(a-d), all demands, including both methods in (c).
\begin{problem}[Changing both bases]\label{ex:04:both-bases}
Let \(T:\R^4\to\R^3\) be
\[
T(x)=(x_1+x_2,\ x_2+x_3-x_4,\ x_1-x_3+x_4)^\top.
\]
Let \(B,C\) be the standard bases of \(\R^4,\R^3\), and let the ordered bases \(E,F\) be
\[
\begin{split}
E={}&\bigl((1,-1,0,-1)^\top,(0,0,1,1)^\top,
             (0,0,0,1)^\top,(1,1,0,0)^\top\bigr),\\
F={}&\bigl((1,0,1)^\top,(1,1,0)^\top,(-1,0,0)^\top\bigr).
\end{split}
\]
\begin{enumerate}
\item Find \([T]_{C\leftarrow B}\).
\item Find both \([I]_{B\leftarrow E}\) and \([I]_{F\leftarrow C}\).
\item Find \([T]_{F\leftarrow E}\) by each of two methods:
directly express the images of the \(E\) basis in \(F\);
then use the change-of-basis formula with the matrices in (a) and (b).
\item Find \(\rank T\) and \(\dim\ker T\).
\end{enumerate}
\end{problem}
\begin{solution}
Write \(F=(f_1,f_2,f_3)\).
\begin{enumerate}
\item The images of the four standard columns give
\[
A=[T]_{C\leftarrow B}
=\begin{pmatrix}1&1&0&0\\0&1&1&-1\\1&0&-1&1\end{pmatrix}.
\]
\item The first matrix has the columns of \(E\):
\[
R=[I]_{B\leftarrow E}
=\begin{pmatrix}1&0&0&1\\-1&0&0&1\\0&1&0&0\\-1&1&1&0\end{pmatrix}.
\]
Its columns are independent: \(Rc=0\) gives
\(c_1+c_4=-c_1+c_4=0\), so \(c_1=c_4=0\);
rows 3 and 4 then give \(c_2=c_3=0\).
For \(y=(y_1,y_2,y_3)^\top\), writing \(y=af_1+bf_2+cf_3\) gives
\[
(a+b-c,b,a)^\top=y,\qquad (a,b,c)=(y_3,y_2,-y_1+y_2+y_3).
\]
Consequently
\[
Q=\begin{pmatrix}1&1&-1\\0&1&0\\1&0&0\end{pmatrix},\qquad
[I]_{F\leftarrow C}=Q^{-1}
=\begin{pmatrix}0&0&1\\0&1&0\\-1&1&1\end{pmatrix}.
\]
These coordinate equations check \(Q^{-1}Q=QQ^{-1}=I_3\).
\item Denote the four members of \(E\) by \(u_1,\ldots,u_4\).
Directly, their images are
\[
T(u_1)=0,\quad T(u_2)=0,\quad
T(u_3)=(0,-1,1)^\top=f_1-f_2,\quad
T(u_4)=(2,1,1)^\top=f_1+f_2,
\]
Thus
\[
H=[T]_{F\leftarrow E}=
\begin{pmatrix}0&0&1&1\\0&0&-1&1\\0&0&0&0\end{pmatrix}.
\]
For the second method, multiplication gives
\[
AR=\begin{pmatrix}0&0&0&2\\0&0&-1&1\\0&0&1&1\end{pmatrix},
\qquad Q^{-1}AR=
\begin{pmatrix}0&0&1&1\\0&0&-1&1\\0&0&0&0\end{pmatrix}=H.
\]
This also verifies \(QH=AR\), the four direct coordinate expansions.
\item In \(H\), columns 3 and 4 are independent:
\(a+b=0\), \(-a+b=0\) imply \(a=b=0\).
All other columns are zero. By the coordinate rank identity following
Theorem~\ref{thm:04:representation},
\(\rank T=\rank H=2\), and Theorem~\ref{thm:04:nullity} gives
\(\dim\ker T=4-2=2\).
\end{enumerate}
\end{solution}

% Harville 22.14; basis extension replaces the cited canonical-form theorem.
\begin{problem}[Bases adapted to a map]\label{ex:04:canonical-map}
Let \(T:V\to W\) be linear of rank \(k>0\), with \(\dim V=n\) and \(\dim W=m\).
Find bases \(E,F\) for which
\[
[T]_{F\leftarrow E}=\begin{pmatrix}I_k&0\\0&0\end{pmatrix}.
\]
\end{problem}
\begin{solution}
\begin{itemize}
\item Choose a basis \(z_1,\ldots,z_{n-k}\) of \(\ker T\) and extend it to a
basis \(E=(x_1,\ldots,x_k,z_1,\ldots,z_{n-k})\) of \(V\).
\item The proof of Theorem~\ref{thm:04:nullity} shows that \(T(x_1),\ldots,T(x_k)\)
form a basis of \(\operatorname{im}T\):
every image is a combination of them, and a zero combination lifts to a
relation between the \(x_i\) and the kernel basis, forcing all coefficients zero.
\item Extend this image basis to a basis \(F\) of \(W\).
The first \(k\) input basis vectors have output coordinate columns
\(e_1,\ldots,e_k\); all remaining input basis vectors map to zero.
These columns give the displayed matrix, of size \(m\times n\).
\end{itemize}
\end{solution}

% Harville 22.16(a,b), following the main textbook's dim(U) in (b).
\begin{problem}[Subspaces preserved by a map]\label{ex:04:invariant-blocks}
Let \(T:V\to V\) be linear, \(\dim V=n\).
A subspace \(U\) is \emph{invariant under \(T\)} when \(T(U)\subseteq U\).
\begin{enumerate}
\item If \(U\) is invariant and \(\dim U=r\), show that there is a basis \(B\) with
\[
[T]_{B\leftarrow B}=\begin{pmatrix}G&H\\0&K\end{pmatrix},
\qquad G\in\R^{r\times r}.
\]
\item If \(U\oplus W=V\) and both subspaces are invariant, find a basis \(B\) with
\[
[T]_{B\leftarrow B}=\operatorname{diag}(G,K),
\qquad G\in\R^{r\times r},\quad r=\dim U.
\]
\end{enumerate}
\end{problem}
\begin{solution}
\begin{enumerate}
\item Extend a basis \(u_1,\ldots,u_r\) of \(U\) to a basis \(B\) of \(V\).
For \(j\leqslant r\), the image \(T(u_j)\) is in \(U\), so its last \(n-r\)
coordinates in \(B\) vanish. These are precisely the lower left block entries.
The remaining entries form the unrestricted blocks \(G,H,K\).
\item Take a basis of \(U\) followed by a basis of \(W\).
The combined list spans \(V\). A zero combination gives an element in
\(U\cap W=\{0\}\), so both separate combinations are zero;
their basis independence proves that the combined list is a basis.
Images of the \(U\) vectors have zero \(W\)-coordinates, and images of
the \(W\) vectors have zero \(U\)-coordinates. Both off-diagonal blocks vanish.
If a subspace is zero, its basis and corresponding blocks are empty,
and the same argument applies.
\end{enumerate}
\end{solution}

\subsection*{Ranks of products and matrix equations}

% MA 4.25(c,d), with proofs avoiding Gram matrices.
\begin{problem}[Multiplication by a full-rank matrix]\label{ex:04:full-rank-products}
For all compatible real matrices, prove
\begin{enumerate}
\item \(\rank(BA)=\rank A\) if \(B\) has full column rank.
\item \(\rank(AC)=\rank A\) if \(C\) has full row rank.
\end{enumerate}
\end{problem}
\begin{solution}
\begin{enumerate}
\item Write \(B:p\times m\), \(A:m\times n\).
The equation \(By=0\) has only \(y=0\), since the columns of \(B\) are independent.
If \(u_1,\ldots,u_r\) is a basis of \(\mathcal C(A)\), then
\(Bu_1,\ldots,Bu_r\) span \(\mathcal C(BA)\).
Moreover, \(\sum_i a_iBu_i=0\) implies \(B(\sum_i a_iu_i)=0\),
then \(\sum_i a_iu_i=0\), and finally all \(a_i=0\).
Thus the two column spaces have equal dimensions.
\item Write \(A:m\times n\), \(C:n\times q\).
Full row rank gives \(\mathcal C(C)=\R^n\).
Every column \(Ax\) in \(\mathcal C(A)\) therefore has
\(x=Cy\) for some \(y\), so \(Ax=ACy\in\mathcal C(AC)\).
The reverse inclusion follows from \(ACy=A(Cy)\).
The column spaces are equal, so their dimensions are equal.
\end{enumerate}
\end{solution}

% Bapat 2.5; no supplied answer. Explicit counterexample supplied here.
\begin{problem}[Removing a middle factor]\label{ex:04:middle-factor}
For \(n\times n\) matrices \(A,B,C\), is it always true that
\(\rank(ABC)\leqslant\rank(AC)\)?
\end{problem}
\begin{solution}
No. Set
\[
A=\begin{pmatrix}1&0\\0&0\end{pmatrix},\quad
B=\begin{pmatrix}0&1\\1&0\end{pmatrix},\quad
C=\begin{pmatrix}0&0\\0&1\end{pmatrix}.
\]
Then \(AC=0\), whereas
\[
BC=\begin{pmatrix}0&1\\0&0\end{pmatrix},\qquad
ABC=\begin{pmatrix}0&1\\0&0\end{pmatrix}.
\]
Thus \(\rank(AC)=0\) and \(\rank(ABC)=1\), disproving the assertion.
\end{solution}

% MA 5.46(c), both original block forms.
\begin{problem}[A zero corner block]\label{ex:04:zero-corner}
Let \(A:m\times n\), \(B:m\times p\), \(C:q\times n\), and \(D:q\times p\).
Prove both inequalities
\[
\rank\begin{pmatrix}A&B\\C&0\end{pmatrix}\geqslant\rank B+\rank C,
\qquad
\rank\begin{pmatrix}0&B\\C&D\end{pmatrix}\geqslant\rank B+\rank C.
\]
\end{problem}
\begin{solution}
Choose independent columns \(b_{j_1},\ldots,b_{j_s}\) forming a basis of
\(\mathcal C(B)\), and columns \(c_{i_1},\ldots,c_{i_r}\) forming a
basis of \(\mathcal C(C)\).
In the first block matrix, select the columns
\[
\begin{pmatrix}b_{j_\ell}\\0\end{pmatrix}\ (1\leqslant\ell\leqslant s),
\qquad
\begin{pmatrix}a_{i_h}\\c_{i_h}\end{pmatrix}\ (1\leqslant h\leqslant r).
\]
In a zero combination, the lower block forces all coefficients of the
second list to vanish. The upper block then forces the coefficients of
the first list to vanish. Thus there are \(r+s\) independent columns.

For the second block matrix, use
\[
\begin{pmatrix}0\\c_{i_h}\end{pmatrix}\ (1\leqslant h\leqslant r),
\qquad
\begin{pmatrix}b_{j_\ell}\\d_{j_\ell}\end{pmatrix}\ (1\leqslant\ell\leqslant s).
\]
Now the upper block first forces the second list's coefficients to vanish;
the lower block then deals with the first list. Again there are \(r+s\)
independent columns. The argument includes empty lists when either rank is zero.
\end{solution}

% MA 5.47(a-c), preserving the specified block argument.
\begin{problem}[A block proof of the product inequalities]\label{ex:04:block-frobenius}
\begin{enumerate}
\item Using Exercise~\ref{ex:04:zero-corner}, prove for
\(A:m\times n\), \(B:n\times p\), \(C:p\times q\) that
\[
\rank(AB)+\rank(BC)\leqslant\rank B+\rank(ABC).
\]
\item Deduce that for \(A:m\times p\), \(B:p\times n\),
\(\rank(AB)\geqslant\rank A+\rank B-p\).
\item With the sizes in (b), deduce that \(AB=0\) implies
\(\rank A\leqslant p-\rank B\).
\end{enumerate}
\end{problem}
\begin{solution}
\begin{enumerate}
\item Put
\[
Z=\begin{pmatrix}0&AB\\BC&B\end{pmatrix},\quad
P=\begin{pmatrix}I_m&-A\\0&I_n\end{pmatrix},\quad
Q=\begin{pmatrix}I_q&0\\-C&I_p\end{pmatrix}.
\]
The matrices obtained from \(P,Q\) by replacing \(-A,-C\) with \(A,C\)
are their inverses: in either multiplication order the off-diagonal
blocks cancel and the diagonal blocks remain identities.
Block multiplication gives
\[
PZ=\begin{pmatrix}-ABC&0\\BC&B\end{pmatrix},\qquad
PZQ=\begin{pmatrix}-ABC&0\\0&B\end{pmatrix}.
\]
For any block diagonal matrix \(\operatorname{diag}(H,K)\),
a basis of its column space is obtained by embedding a basis of
\(\mathcal C(H)\) in the upper coordinates and a basis of
\(\mathcal C(K)\) in the lower coordinates.
These columns span; a zero combination, read in the two blocks separately,
has every coefficient zero. Therefore its rank is \(\rank H+\rank K\).
Invertible multiplication preserves rank, and multiplication by \(-1\)
preserves a column space. Hence
\[
\rank Z=\rank(ABC)+\rank B.
\]
Exercise~\ref{ex:04:zero-corner} also gives
\(\rank Z\geqslant\rank(AB)+\rank(BC)\), proving (a).
\item Apply (a) to the three factors \(A,I_p,B\):
\[
\rank A+\rank B\leqslant p+\rank(AB).
\]
\item In (b), substitute \(\rank(AB)=0\) and rearrange.
\end{enumerate}
\end{solution}

% Bapat 2.17; expand the rank-factorization hint.
\begin{problem}[Using the rank of the middle factor]\label{ex:04:factor-frobenius}
Let \(A,X,B\in\R^{n\times n}\). Prove
\[
\rank(AXB)\geqslant\rank(AX)+\rank(XB)-\rank X
\]
by using a rank factorization of \(X\).
\end{problem}
\begin{solution}
\begin{itemize}
\item If \(X=0\), every term is zero. Otherwise let \(r=\rank X>0\)
and choose \(X=UV\), with \(U:n\times r\) of full column rank
and \(V:r\times n\) of full row rank.
\item Exercise~\ref{ex:04:full-rank-products} gives
\[
\rank(AX)=\rank(AUV)=\rank(AU),\qquad
\rank(XB)=\rank(UVB)=\rank(VB).
\]
\item The factors \(AU:n\times r\) and \(VB:r\times n\) have inner size \(r\).
Applying Proposition~\ref{prop:04:bounds} to them gives
\[
\rank(AXB)=\rank((AU)(VB))
\geqslant\rank(AU)+\rank(VB)-r,
\]
which is the required inequality.
\end{itemize}
\end{solution}

% Harville 7.1(a,b), all four assertions.
\begin{problem}[When cancellation is valid]\label{ex:04:cancellation}
\begin{enumerate}
\item Let \(A:m\times n\), \(C:n\times q\), and \(B:q\times p\).
Prove both statements:
\[
\begin{split}
\rank(AC)=\rank C&\ \Longrightarrow\
\mathcal R(ACB)=\mathcal R(CB),\quad \rank(ACB)=\rank(CB),\\
\rank(CB)=\rank C&\ \Longrightarrow\
\mathcal C(ACB)=\mathcal C(AC),\quad \rank(ACB)=\rank(AC).
\end{split}
\]
\item Let \(A,B:m\times n\).
Prove that
\[
\begin{split}
C:r\times q,\ D:q\times m,\ \rank(CD)=\rank D,\ CDA=CDB
&\ \Longrightarrow\ DA=DB,\\
C:n\times q,\ D:q\times p,\ \rank(CD)=\rank C,\ ACD=BCD
&\ \Longrightarrow\ AC=BC.
\end{split}
\]
\end{enumerate}
\end{problem}
\begin{solution}
\begin{enumerate}
\item Rows of a product are combinations of the rows of its right factor.
Thus \(\mathcal R(AC)\subseteq\mathcal R(C)\).
Under the first rank hypothesis these spaces have equal dimensions, hence are equal.
Each row of \(C\) can therefore be expressed using the rows of \(AC\);
collecting those coefficients gives a matrix \(L:n\times m\) with \(C=LAC\).
Now
\[
CB=LACB,\qquad ACB=A(CB),
\]
so each of \(\mathcal R(CB)\) and \(\mathcal R(ACB)\) is contained in the other.
Their equality gives equality of ranks.

Under the second hypothesis, \(\mathcal C(CB)\subseteq\mathcal C(C)\)
and equal dimensions give equality.
Express the columns of \(C\) using the columns of \(CB\):
\(C=CBR\) for some \(R:p\times q\).
The identities
\[
AC=ACBR,\qquad ACB=(AC)B
\]
give both inclusions of column spaces and hence both claimed equalities.
\item Put \(F=A-B\).
For the first assertion, the first rank statement in (a), applied
to the factors \(C,D,F\), gives
\[
\rank(DF)=\rank(CDF)=0,
\]
since \(CDF=CDA-CDB=0\). Therefore \(DF=0\), or \(DA=DB\).
For the second assertion, use the second rank statement in (a)
with factors \(F,C,D\):
\[
\rank(FC)=\rank(FCD)=0.
\]
Thus \(FC=0\), giving \(AC=BC\).
\end{enumerate}
\end{solution}

% MA 6.47, original two-parameter system; no determinant argument.
\begin{problem}[Parameters and consistency]\label{ex:04:parameters}
For real \(\alpha,\beta\), consider
\[
\begin{aligned}
\alpha x_1+\beta x_2+2x_3&=1,\\
\alpha x_1+(2\beta-1)x_2+3x_3&=1,\\
\alpha x_1+\beta x_2+(\beta+3)x_3&=2\beta-1.
\end{aligned}
\]
When is the system consistent? When is its solution unique?
\end{problem}
\begin{solution}
Subtract the first equation from the second and from the third.
Then replace the first by its difference with the new second equation.
These reversible operations give
\[
\begin{aligned}
\alpha x_1+x_2+x_3&=1,\\
(\beta-1)x_2+x_3&=0,\\
(\beta+1)x_3&=2(\beta-1).
\end{aligned}
\]
No parameter has been divided out.
\begin{itemize}
\item If \(\beta=-1\), the last equation is \(0=-4\), so there is no solution.
\item If \(\beta=1\), the last two equations give \(x_3=0\), and the first gives
\(x_2=1-\alpha x_1\). Thus
\[
(x_1,x_2,x_3)=(t,1-\alpha t,0)\qquad(t\in\R),
\]
so there are infinitely many solutions for every \(\alpha\).
\item Suppose \(\beta\ne-1,1\). Back substitution gives
\[
x_3=\frac{2(\beta-1)}{\beta+1},\qquad
x_2=-\frac{2}{\beta+1},\qquad
\alpha x_1=\frac{5-\beta}{\beta+1}.
\]
If \(\alpha\ne0\), this determines the unique solution
\[
(x_1,x_2,x_3)=
\left(\frac{5-\beta}{\alpha(\beta+1)},
      -\frac{2}{\beta+1},\frac{2(\beta-1)}{\beta+1}\right).
\]
If \(\alpha=0\), consistency requires \(\beta=5\).
For \(\alpha=0,\beta=5\), all solutions are
\((t,-1/3,4/3)\), \(t\in\R\); for the other values in this case there is no solution.
\end{itemize}
The displayed families satisfy the transformed equations, so reversing
the row operations verifies them in the original system.
Thus uniqueness holds exactly when \(\alpha\ne0\) and \(\beta\ne-1,1\).
\end{solution}
